\subsection{Roots of unity}
Roots of unity are an important concept in Galois-theory. Galois-theory 
takes advantages of some key results made by Ernst Eduard Kummer around 1840s. Kummer 
studied the properties of roots of unity, especially in perfect fields. Galois later 
discovered a profound relation with the so called radical extensions. 

\begin{definition}[Perfect field]
    A perfect field, is a field with no two non-zero elements such that their 
    product is zero. Such field is sometimes also sometimes called a zero-characteristic
    field or an integral field (rings are also sometimes called integral).  
\end{definition}
\vspace{0.2em}

\begin{definition}[Root of unity]
    In a field, a root of unity of order $k$ is an element $\omega$ such that 
    $\omega^k = 1$. Of course, $1$ is always root of unity and it's called 
    trivial.  
\end{definition}
\vspace{0.2em}

\begin{lemma}
\label{lem:RootsOfUnityCyclic}
    In every field, there is a so called primitive root of unity $\omega_i$ such that every 
    other non-trivial roots of unity $\omega_i$ of the same order can be expressed as 
    a power of $\omega$. 
\end{lemma}
\begin{proof}
    WORK IN PROGRESS.
\end{proof}    
\vspace{0.2em}

\begin{lemma}
\label{lem:FOmegaSolvable}
    The Galois group $Gal_{F}(F(\omega))$ for some field $F$ and some 
    non-trivial root of unity $\omega$ is abelian, herby it's solvable. 
\end{lemma}
\begin{proof}
    Consider the product (composition) of two automorphisms $\varphi$, $\psi$ 
    randomly choosen in $Gal_{F}(F(\omega))$. They commute:

    $$ \psi(\varphi(\omega^i)) = \psi(\omega^{ij}) = \omega^{ijk} = 
    \omega^{ikj} = \varphi(\omega^{ik}) = \varphi(\psi(\omega^i)) $$
\end{proof}

\newpage